\documentclass[10pt,letterpaper]{article}
\usepackage[T1]{fontenc}
\usepackage[utf8]{inputenc}
\usepackage[sfdefault]{sourcesanspro}
\usepackage{inconsolata}
\usepackage[margin=0.84in,headheight=15pt]{geometry}
\usepackage{microtype}
\usepackage{xcolor}
\usepackage{booktabs,longtable,tabularx,array}
\usepackage{enumitem}
\usepackage{fancyhdr,lastpage}
\usepackage[most]{tcolorbox}
\usepackage{listings}
\usepackage{titlesec}
\usepackage{hyperref}
\usepackage{amsmath,amssymb}
\usepackage{multicol}
\definecolor{navy}{HTML}{15324C}
\definecolor{teal}{HTML}{087E83}
\definecolor{slate}{HTML}{536879}
\definecolor{paper}{HTML}{F2F6F8}
\definecolor{warm}{HTML}{E8F4F3}
\hypersetup{colorlinks=true,urlcolor=teal,linkcolor=navy,pdfauthor={Software Engineering Reference Suite},pdfsubject={MISRA C++:2023 engineering guidance}}
\titleformat{\section}{\Large\bfseries\color{navy}}{\thesection}{0.7em}{}
\titleformat{\subsection}{\large\bfseries\color{teal}}{\thesubsection}{0.6em}{}
\titleformat{\subsubsection}{\normalsize\bfseries\color{navy}}{\thesubsubsection}{0.5em}{}
\setlist[itemize]{leftmargin=*,topsep=2pt,itemsep=2pt}
\setlist[enumerate]{leftmargin=*,topsep=2pt,itemsep=3pt}
\setlength{\parskip}{5pt}
\setlength{\parindent}{0pt}
\setlength{\emergencystretch}{2em}
\pagestyle{fancy}
\fancyhf{}
\fancyhead[L]{\small\color{slate}MISRA C++:2023 Engineering Suite}
\fancyhead[R]{\small\color{slate}\nouppercase{\leftmark}}
\fancyfoot[L]{\scriptsize\color{slate}Independent engineering guidance; see licensed source}
\fancyfoot[R]{\scriptsize\color{slate}\thepage\ / \pageref{LastPage}}
\newtcolorbox{sourcebox}{colback=paper,colframe=navy!30,boxrule=0.45pt,arc=1.6pt,breakable}
\newtcolorbox{practicebox}{colback=warm,colframe=teal!40,boxrule=0.45pt,arc=1.6pt,breakable}
\newcommand{\SourceNote}[1]{\begin{sourcebox}\textbf{SOURCE BASIS / NORMATIVE POINTER}\quad #1\end{sourcebox}}
\newcommand{\PracticeNote}[1]{\begin{practicebox}\textbf{ENGINEERING RECOMMENDATION / NOT MISRA TEXT}\quad #1\end{practicebox}}
\newcommand{\Record}[1]{\textsf{\bfseries #1}}
\lstdefinestyle{cppsuite}{language=C++,basicstyle=\ttfamily\footnotesize,keywordstyle=\color{navy}\bfseries,commentstyle=\color{slate},stringstyle=\color{teal},breaklines=true,breakatwhitespace=true,columns=fullflexible,frame=single,framerule=0.3pt,rulecolor=\color{navy!35},backgroundcolor=\color{paper},showstringspaces=false,tabsize=2,xleftmargin=0.5em,xrightmargin=0.5em,aboveskip=4pt,belowskip=4pt}
\lstset{style=cppsuite}
\newcommand{\SuiteTitle}[3]{%
\begin{titlepage}\thispagestyle{empty}
\vspace*{1.5cm}
{\color{teal}\rule{2.4cm}{2pt}}\par\vspace{1cm}
{\Huge\bfseries\color{navy}MISRA C++:2023\\[0.4em]#1}\par\vspace{0.7cm}
{\Large\color{slate}#2}\par\vspace{1.2cm}
\begin{sourcebox}\textbf{Source baseline:} MISRA C++:2023, Guidelines for the use of C++17 in critical systems, October 2023 (user-supplied licensed copy). This is an independent companion, not an official MISRA publication, rule replacement, certified checker, or claim of project compliance.\end{sourcebox}
\vfill{\large\color{navy}Technical edition \textbullet\ 16 October 2026}\par\smallskip
{\small\color{slate}#3}\par\end{titlepage}
\tableofcontents\clearpage
}
\newcommand{\StandardsBoundary}{%
\section*{How to read normative status}\addcontentsline{toc}{section}{How to read normative status}
\textbf{Normative source:} the licensed MISRA C++:2023 guideline text, including any applicable \emph{Amplification} and \emph{Exception}. A guideline \emph{headline} alone is not a complete compliance test. The source's Section~3 defines classifications, categories, decidability, analysis scope, and applicability.\par
\textbf{Independent guidance:} all review questions, programming patterns, lifecycle processes, sample thresholds, ownership assignments, controls, and CI configurations in this suite are \emph{engineering recommendations}, not statements issued by MISRA.\par
\textbf{Critical limitation:} no individual example is asserted compliant with the whole standard. For formal compliance claims, follow the applicable edition of \emph{MISRA Compliance} as the source itself instructs, with a documented compliance plan and supporting evidence.\par
}

\begin{document}
\SuiteTitle{Training Curriculum}{12-week dependency-ordered course with assessments and labs}{Volume 11 / Professional Development}
\StandardsBoundary

\section{Curriculum policy}
The 12-week syllabus progresses from C++17 language semantics to source interpretation, static-analysis methods and compliance evidence. It assumes participants can write and debug modern C++, and it requires access to the licensed MISRA C++:2023 document for correct interpretation. Learning hours are suggested, not prescribed by MISRA.

The dependency structure is cumulative: language foundations precede rule analysis; rule analysis precedes automated adjudication; adjudication precedes release claims.


\section{Week-by-week learning plan}
\small\setlength{\tabcolsep}{3pt}\renewcommand{\arraystretch}{1.15}
\begin{longtable}{p{0.07\linewidth}p{0.23\linewidth}p{0.14\linewidth}p{0.46\linewidth}}
\toprule
Week & Module & Prerequisite & Objective/primary output\\ \midrule
\endfirsthead
\toprule
Week & Module & Prerequisite & Objective/primary output\\ \midrule
\endhead
1 & Environment and source taxonomy & C++17 basics & Correctly name directive/rule, Mandatory/Required/Advisory, applicability, amplification and exception. Output: Short quiz and baseline manifest.\\
2 & Critical C++ language semantics & Week 1 & Reason about undefined/unspecified behavior, integer widths, initialization and sequencing. Output: Memory and semantics worksheet.\\
3 & Types and conversions & Week 2 & Identify conversion chains, casts, null pointers and explicit predicates. Output: Type-transition worksheet and fixture review.\\
4 & Control flow and value use & Weeks 2--3 & Draw CFG, call graphs, data-flow definitions and unreachable locations. Output: CFG + defect reproductions.\\
5 & Declarations and object model & Weeks 2--4 & Reason about ODR, linkage, lifetime, initialization and class state. Output: Two-TU analysis exercise.\\
6 & Inheritance, templates, operators & Week 5 & Analyze overrides, dynamic dispatch, copy/move semantics and template usage. Output: Class review dossier.\\
7 & Exception and preprocessor design & Weeks 3--6 & Review throw/catch policies, macros, includes and build variants. Output: Exception graph and preprocessing report.\\
8 & STL, memory, and I/O & Weeks 5--7 & Evaluate ownership, containers, algorithms, C library facilities and I/O. Output: RAII refactor and invalidation checklist.\\
9 & Rule-by-rule inspection practicum & Weeks 1--8 & Apply all source guidelines to topical specimen clusters with source amplification checks. Output: 179-ID worksheet with source cross-references.\\
10 & Static analysis and quality evidence & Weeks 4--9 & Measure D/U and TU/System tool coverage and adjudicate warnings. Output: Analyzer limitations and coverage report.\\
11 & Compliance and deviation governance & Weeks 9--10 & Apply valid category-specific dispositions and independent approval. Output: Mock compliance plan + deviation request.\\
12 & CI/CD and capstone assurance & Weeks 1--11 & Prove build context and create an evidence-backed, scoped assessment. Output: Reviewed capstone evidence dossier.\\
\bottomrule
\end{longtable}
\normalsize

\section{Reading and practice hour model}
\small\setlength{\tabcolsep}{3pt}\renewcommand{\arraystretch}{1.15}
\begin{longtable}{p{0.09\linewidth}p{0.21\linewidth}p{0.58\linewidth}}
\toprule
Week & Suggested effort & Required reference\\ \midrule
\endfirsthead
\toprule
Week & Suggested effort & Required reference\\ \midrule
\endhead
1 & 2h study + 3h lab & Source Sections 1--3\\
2 & 3h study + 3h lab & Sections 4.4, 4.6\\
3 & 3h study + 3h lab & Sections 4.7--4.8\\
4 & 3h study + 3h lab & Sections 4.0 and 4.9\\
5 & 4h study + 3h lab & Sections 4.6, 4.10--4.12\\
6 & 3h study + 4h lab & Sections 4.13--4.17\\
7 & 3h study + 3h lab & Sections 4.18--4.19\\
8 & 4h study + 3h lab & Sections 4.21--4.30\\
9 & 4h study + 5h lab & Volumes 02--05; source Chapter 4\\
10 & 3h study + 4h lab & Source Sections 3.5--3.7; Vol 06\\
11 & 4h study + 4h lab & Source Section 1.5, Sections 3.1--3.8; MISRA Compliance separately\\
12 & 3h study + 7h lab & Volumes 07--10\\
\bottomrule
\end{longtable}
\normalsize

\section{Knowledge-to-evidence competency map}
\small\setlength{\tabcolsep}{3pt}\renewcommand{\arraystretch}{1.15}
\begin{longtable}{p{0.2\linewidth}p{0.34\linewidth}p{0.34\linewidth}}
\toprule
Competency & Knowledge evidence & Demonstration\\ \midrule
\endfirsthead
\toprule
Competency & Knowledge evidence & Demonstration\\ \midrule
\endhead
C++17 semantics & Explain object lifetime and conversion/sequence rules. & Correct live code with tests and source citations.\\
Rule interpretation & Find normative amplification, exception, scope, category. & Complete case study with source-referenced verdict.\\
Automated analysis & Identify checker accuracy and interprocedural limits. & Tool mapping/uncertainty analysis on two build variants.\\
Compliance governance & Apply Mandatory/Required/Advisory dispositions. & Review deviation record and mock compliance claim.\\
Delivery engineering & Relate scanner execution to shipped artifact. & Reproducible CI audit bundle with independent acceptance.\\
\bottomrule
\end{longtable}
\normalsize

\section{Assessment strategy}
\begin{itemize}
\item Weekly checks: 10--15 scenario questions, one code correction, one source-location exercise.
\item Milestone A (week 4): CFG and types assessment; recommended pass mark 80%.
\item Milestone B (week 8): class/lifetime/standard-library technical review; recommended pass mark 85%.
\item Milestone C (week 11): category and deviation application; no incorrect approval of Mandatory violations.
\item Final capstone (week 12): 85% weighted score, documented build fidelity, complete evidence scope, reviewer signoff.
\item Retraining triggers: new guideline edition, compiler/target changes, tool check-map changes or demonstrated review error.
\end{itemize}

\section{Prerequisite graph}
\begin{lstlisting}[language={},caption={Directed dependency ordering; each line lists direct prerequisites}]
W1: none
W2: W1
W3: W2
W4: W2, W3
W5: W2, W4
W6: W5
W7: W3, W6
W8: W5, W7
W9: W1, W2, W3, W4, W5, W6, W7, W8
W10: W4, W9
W11: W9, W10
W12: W10, W11
\end{lstlisting}

\section{Instructor controls}
For every assessment, sample solutions should identify which statements are derived from the licensed source and which are engineering interpretation. Instructors must not convert a warning heuristic into a new MISRA requirement. Maintain a dated source/version register and requalify materials when official clarifications affect the exercises.


\end{document}
